\subsection{Barrels and barrelled spaces}

PROPOSITION 1. --- Let T be a subset of E. The following conditions are equivalent :

\begin{enumerate}
\def\labelenumi{(\roman{enumi})}
\item
  T is convex, balanced, closed and absorbent.
\item
  T is the polar of a convex, balanced and weakly bounded set M in E'.
\item
  There exists a lower semi-continuous semi-norm \(p\) on \(\mathbf{E}\) , such that \(\mathbf{T}\) is the set of all \(x \in \mathbf{E}\) satisfying \(p(x) \leqslant 1\) .
\item
  \(\Rightarrow\) (ii): under the hypothesis of (i), let \(\mathbf{M} = \mathbf{T}^{\circ}\) ; then \(\mathbf{M}\) is convex and balanced in \(\mathbf{E}'\) . For every \(x \in \mathbf{E}\) , there exists a real number \(r > 0\) such that \(rx \in \mathbf{T}\) , hence \(|u(x)| \leqslant \frac{1}{r}\) for all \(u \in \mathbf{M}\) ; in other words \(\mathbf{M}\) is weakly bounded. From cor. 3 of II, p.~45, we have \(\mathrm{T} = \mathbf{M}^{\circ}\) , hence \(\mathrm{T}\) satisfies (ii).
\item
  \(\Rightarrow\) (iii): under the hypothesis of (ii), let \(p(x) = \sup_{u\in \mathbf{M}}|u(x)|\) for all \(x\in \mathrm{E}\) . It is immediate that \(T = M^{\circ}\) consists of all \(x \in E\) such that \(p(x) \leqslant 1\) . The semi-norm p on \(E'\) is lower semi-continuous, being the superior envelope of a family of continuous functions (GT, IV, § 6, No.~2, corollary).
\item
  \(\Rightarrow\) (i): this is clear.
\end{enumerate}

COROLLARY. --- The following conditions are equivalent :

\begin{enumerate}
\def\labelenumi{(\roman{enumi})}
\item
  every weakly bounded subset of \(\mathrm{E}'\) is equicontinuous;
\item
  every convex, balanced, closed and absorbent set in E is a neighbourhood of 0;
\item
  every lower semi-continuous semi-norm on E is continuous.
\end{enumerate}

DEFINITION 1. --- A set T satisfying the equivalent conditions of prop. 1 is said to be a barrel in E.

DEFINITION 2. --- A space E is said to be barrelled if it satisfies the equivalent conditions of the corollary of prop. 1.

We know (III, p.~22, prop. 9) that every equicontinuous subset of the dual E' of E is strongly and weakly bounded. We can therefore restate the definition of barrelled spaces as follows :

Scholium. --- In the dual of a barrelled space, the class of equicontinuous sets, the class of strongly bounded sets, the class of weakly bounded sets and the class of relatively compact sets for the weak topology are all identical. If E is Hausdorff and barrelled, and if \(E_{b}'\) is its strong dual, the polars of the neighbourhoods of 0 in one of the spaces form a base of the canonical bornology of the other, and the polars of bounded subsets of one of the spaces form a base for the filter of neighbourhoods of 0 of the other space.

PROPOSITION 2. --- Every locally convex space E which is a Baire space (GT, IX, § 5, No.~3) is barrelled.

Let \(\mathrm{T}\) be a barrel in \(\mathrm{E}\) ; since \(\mathrm{T}\) is absorbent, \(\mathrm{E}\) is the union of closed sets \(n\mathrm{T}\) ( \(n\) integer \(> 0\) ); since \(\mathrm{E}\) is a Baire space, at least one of these sets contains an interior point, hence T itself has an interior point x. If \(x \neq 0\) , since \(-x \in T\) , and 0 is a point of the open segment with extremities x and -x, 0 is an interior point of the convex set T (II, p.~14, prop. 16). Therefore T is a neighbourhood of 0.

COROLLARY. --- Every Fréchet space (and in particular, every Banach space) is barrelled. This follows from Baire's theorem (GT, IX, § 5, No.~3, th. 1).

PROPOSITION 3. --- Let \((\mathrm{F}_{i})_{i\in\mathrm{I}}\) be a family of barrelled spaces, and for every \(i\in I\) , let \(f_{i}\) be a linear mapping from \(F_{i}\) into a vector space E. The space E with the finest locally convex topology for which the \(f_{i}\) are continuous (II, p.~27, prop. 5), is a barrelled space.

Let T be a barrel in E. Since \(f_{i}\) is continuous, \(f_{i}^{-1}(T)\) is a convex, balanced, closed and absorbent set in \(F_{i}\) ; in other words, a barrel in \(F_{i}\) ; since \(F_{i}\) is barrelled, \(f_{i}^{-1}(T)\) is a neighbourhood of 0 in \(F_{i}\) , for all \(i \in I\) . This implies that T is a neighbourhood of 0 in E (II, p.~27, prop. 5).

COROLLARY 1. --- Every quotient space of a barrelled space is barrelled.

COROLLARY 2. --- Let \((\mathrm{E}_{i})_{i\in\mathrm{I}}\) be a family of locally convex spaces and let E be the topological direct sum of this family. For E to be barrelled, it is necessary and sufficient that each \(E_{i}\) is barrelled.

The condition is evidently sufficient by virtue of prop. 3; it is also necessary, by cor. 1, since each \(E_{i}\) is isomorphic to a quotient space of E (II, p.~31, prop. 8).

COROLLARY 3. --- Every inductive limit of barrelled spaces is a barrelled space.

We shall prove later (IV, p.~14, corollary) that every product of barrelled spaces is barrelled.

